\section*{Capítulo \ref{ch:b2:diffcalc} --- Cálculo diferencial}

\begin{solution}{exo:b2:diffcalc:1}
$J_f = \begin{pmatrix} 2x & -2y\\ 2y & 2x\end{pmatrix}$, $\det J_f =
4(x^2 + y^2)$: invertible fuera del origen. (Esta $f$ es $z \mapsto
z^2$ disfrazada de real.)

$J_\Phi = \begin{pmatrix} \cos\theta & -r\sin\theta & 0\\ \sin\theta
& r\cos\theta & 0\\ 0 & 0 & 1\end{pmatrix}$, $\det = r$: invertible
para $r \neq 0$ (coordenadas cilíndricas).
\end{solution}

\begin{solution}{exo:b2:diffcalc:2}
Para $v = (a, b) \neq 0$: $\frac{f(tv) - 0}{t} = \frac{t^3a^3}{t\cdot
t^2(a^2+b^2)} = \frac{a^3}{a^2 + b^2}$: existe toda derivada
direccional, con valor $D_v = \frac{a^3}{a^2+b^2}$. Pero $v \mapsto
D_v$ no es lineal ($D_{(1,0)} = 1$, $D_{(0,1)} = 0$, $D_{(1,1)} =
\frac12 \neq 1$): ninguna aplicación lineal puede producir esos
valores, luego $f$ no es \omterm{def:b2:diffcalc:differential}{diferenciable} en $0$ (la \omterm{def:b2:diffcalc:differential}{diferencial}
tendría que ser $v \mapsto D_v$).
\end{solution}

\begin{solution}{exo:b2:diffcalc:3}
$f = x^3 + y^3 - 3xy$: puntos críticos $(0,0)$ y $(1,1)$ (cálculo del
primer año). Hessianas: $H = \begin{pmatrix} 6x & -3\\ -3 &
6y\end{pmatrix}$. En $(0,0)$: signos propios mixtos ($\det = -9 <
0$): punto de silla. En $(1,1)$: $\det = 27 > 0$ y traza $> 0$:
definida positiva, mínimo local estricto; ahora justificado por el~\cref{thm:b2:diffcalc:extrema} en vez de decretado.

$g = x^4 + y^4 - 2(x-y)^2$: $\nabla g = (4x^3 - 4(x - y),\; 4y^3 +
4(x-y))$; puntos críticos $(0,0)$, $(\sqrt2, -\sqrt2)$ y $(-\sqrt2,
\sqrt2)$ (primer año). En $(\pm\sqrt2, \mp\sqrt2)$: $H =
\begin{pmatrix} 12\cdot2 - 4 & 4\\ 4 & 20\end{pmatrix} =
\begin{pmatrix} 20 & 4\\ 4 & 20\end{pmatrix}$: definida positiva
(diagonalmente dominante; \omterm{def:b2:reduction:eigen}{valores propios} $24, 16$): mínimos locales
estrictos. En $(0,0)$: $H = \begin{pmatrix} -4 & 4\\ 4 &
-4\end{pmatrix}$, singular y semidefinida negativa: el test calla; el
estudio direccional ($g(x,x) = 2x^4 > 0$, $g(x,-x) = 2x^4 - 8x^2 < 0$
para $x$ pequeño) muestra un punto de tipo silla: no hay extremo.
\end{solution}

\begin{solution}{exo:b2:diffcalc:4}
Derívese $t \mapsto f(tx)$ en $t = 1$: por la regla de la cadena,
$\langle \nabla f(x), x\rangle$; y por homogeneidad, esa misma
función es $t^pf(x)$, de derivada $pf(x)$ en $t = 1$: la identidad de
Euler.

Recíproco: fíjese $x \neq 0$ y sea $\varphi(t) = f(tx) - t^p f(x)$
sobre $t > 0$. Entonces $\varphi'(t) = \langle\nabla f(tx), x\rangle
- pt^{p-1}f(x) = \frac1t\bigl(\langle \nabla f(tx), tx\rangle -
p\,t^pf(x)\bigr)$. La hipótesis ---la identidad de Euler en el punto
$tx$--- evalúa el corchete como $p\,f(tx) - p\,t^pf(x) =
p\,\varphi(t)$. Así pues, $\varphi' = \frac{p}{t}\varphi$ con
$\varphi(1) = 0$: la única solución de esa EDO lineal es $\varphi
\equiv 0$ (unicidad del primer año), es decir, $f(tx) = t^pf(x)$.
\end{solution}

\begin{solution}{exo:b2:diffcalc:5}
Desarrollando $f(x + h) - f(x) = \langle Ax - b, h\rangle +
\frac12\langle Ah, h\rangle$ (simetría de $A$): $\nabla f(x) = Ax -
b$ y $H_f = A$ en todas partes. $f$ es convexa si y solo si $A$ es
semidefinida positiva (el test de la hessiana, aquí global porque $H$
es constante: la fórmula de Taylor de segundo orden es exacta). Si
$A$ es definida positiva, el único punto crítico es $x^* = A^{-1}b$,
y $f(x^* + h) - f(x^*) = \frac12\langle Ah, h\rangle \geq
\frac{\lambda_{\min}}{2}\norm h^2 > 0$ para $h \neq 0$: mínimo global
estricto.
\end{solution}

\begin{solution}{exo:b2:diffcalc:6}
Como $\nabla g(a) \neq 0$, alguna parcial, digamos $\frac{\partial
g}{\partial y}(a) \neq 0$: el teorema de la función implícita (el
compañero del~\cref{thm:b2:diffcalc:inverse}) parametriza $\{g = 0\}$
cerca de $a = (a_1, a_2)$ como $\gamma(t) = (t, y(t))$ con $y$ de
clase $C^1$ y $y'(t) = -\frac{\partial_x g}{\partial_y
g}(\gamma(t))$ (derívese $g(t, y(t)) = 0$). La función de una
variable $t \mapsto f(\gamma(t))$ tiene un extremo local en $t =
a_1$:
\[
0 = \frac{\dd}{\dd t}f(\gamma(t))\Big|_{a_1}
= \partial_x f(a) + \partial_y f(a)\,y'(a_1)
= \partial_x f(a) - \partial_yf(a)\frac{\partial_x
g(a)}{\partial_y g(a)} :
\]
los vectores $\nabla f(a)$ y $\nabla g(a)$ tienen coordenadas
proporcionales: $\nabla f(a) = \lambda \nabla g(a)$ con $\lambda =
\frac{\partial_y f(a)}{\partial_y g(a)}$.

Aplicación: sobre la circunferencia, que $\nabla(xy) = (y, x)$ sea
paralelo a $(2x, 2y)$ fuerza $y^2 = x^2$; con la restricción, los
candidatos son $\pm\bigl(\tfrac{1}{\sqrt2},
\tfrac{1}{\sqrt2}\bigr)$ (de valor $\frac12$) y
$\pm\bigl(\tfrac{1}{\sqrt2}, -\tfrac{1}{\sqrt2}\bigr)$ (de valor
$-\frac12$): máximo $\frac12$ y mínimo $-\frac12$ (alcanzados: la
circunferencia es \omterm{def:b2:metric:compact}{compacta}).
\end{solution}

\begin{solution}{exo:b2:diffcalc:7}
$\nabla f = (2x - y + 1,\; 2y - x - 1) = 0$: resolviendo, $x =
-\frac13$, $y = \frac13$. Hessiana $\begin{pmatrix} 2 & -1\\ -1 &
2\end{pmatrix}$, definida positiva: mínimo global de la $f$
cuadrática, de valor $f\bigl(-\frac13, \frac13\bigr) = -\frac13$.

Sobre el triángulo $T$: el punto crítico interior $(-\frac13,
\frac13) \notin T$ ($x$ negativo). Lados: sobre $y = 0$, $x \in
\intcc{0}{1}$: $f = x^2 + x$, creciente, con extremos $0$ y $2$.
Sobre $x = 0$: $f = y^2 - y$, con mínimo $-\frac14$ en $y =
\frac12$, y valores $0$ y $0$ en los extremos. Sobre $x + y = 1$:
sustituyendo $y = 1 - x$, $f = x^2 + (1-x)^2 - x(1-x) + x - (1-x) =
3x^2 - x$; sobre $\intcc{0}{1}$: mínimo $-\frac{1}{12}$ en $x =
\frac16$, y valores $0$ (en $x=0$) y $2$ (en $x=1$). Vértices:
$f(0,0) = 0$, $f(1,0) = 2$, $f(0,1) = 0$. Globalmente sobre $T$:
mínimo $-\frac14$ en $(0, \frac12)$ y máximo $2$ en $(1, 0)$.
\end{solution}

\begin{solution}{exo:b2:diffcalc:8}
$J_f = \begin{pmatrix} 1 & 2y\\ 2x & 1\end{pmatrix}$, $\det J_f = 1 -
4xy$. En $0$: $\det = 1 \neq 0$, luego difeomorfismo local
(\cref{thm:b2:diffcalc:inverse}), con
\[
\dd(f^{-1})_{(0,0)} = (J_f(0))^{-1} = I_2 .
\]
Invertibilidad sobre una bola: hace falta $4\abs{xy} < 1$ en todo
punto; sobre $\norm{(x,y)}_2 < r$, $\abs{xy} \leq \frac{x^2 +
y^2}{2} < \frac{r^2}{2}$, de modo que $r = \frac{1}{\sqrt2}$ sirve; y
es el mayor posible: en $(x, y) = \bigl(\tfrac12, \tfrac12\bigr)$,
de \omterm{def:b2:nvs:norm}{norma} $\frac{1}{\sqrt2}$, se tiene $\det J_f = 0$.
\end{solution}

\begin{solution}{exo:b2:diffcalc:9}
$f(t) = (\cos t, \sin t)$: $f(0) = f(2\pi) = (1, 0)$, y sin embargo
$f'(t) = (-\sin t, \cos t)$ tiene \omterm{def:b2:nvs:norm}{norma} $1$ y nunca se anula: no hay
ningún punto donde la derivada se anule; Rolle no tiene análogo
vectorial. La desigualdad del valor medio se cumple holgadamente:
$\norm{f(2\pi) - f(0)} = 0 \leq 2\pi \cdot \sup\norm{f'} = 2\pi$.
\end{solution}

\begin{solution}{exo:b2:diffcalc:10}
$f(x+h) - f(x) = 2\langle x, h\rangle + \norm h^2$: $\dd f_x =
2\langle x, \cdot\rangle$, $\nabla f(x) = 2x$ y hessiana $2I$
(constante). Para $g$:
\[
g(x + h) - g(x) = \langle Ax, h\rangle + \langle Ah, x\rangle
+ \langle Ah, h\rangle
= \bigl\langle (A + A^{\mathsf T})x,\ h\bigr\rangle +
O(\norm h^2),
\]
luego $\nabla g(x) = (A + A^{\mathsf T})x$ y $H_g = A + A^{\mathsf
T}$, constante. Por el~\cref{exo:b2:diffcalc:11}, $g$ es convexa si y
solo si $A + A^{\mathsf T}$ es semidefinida positiva: solo importa la
parte \omterm{def:b2:quadratic:adjoint}{simétrica} de $A$, como confirma que $g(x) = \langle \frac{A +
A^{\mathsf T}}2 x, x\rangle$.
\end{solution}

\begin{solution}{exo:b2:diffcalc:11}
$f$ es convexa si y solo si su restricción a todo segmento es
convexa, es decir, si y solo si toda $\varphi(t) = f(a + tv)$ es
convexa. Por la regla de la cadena, $\varphi''(t) = \langle
H_{a+tv}\,v,\ v\rangle$.

Si todas las hessianas son semidefinidas positivas, $\varphi'' \geq
0$, luego cada $\varphi$ es convexa (volumen del primer año) y $f$ es
convexa. Recíprocamente, si $f$ es convexa, cada $\varphi$ lo es,
luego $\varphi''(0) \geq 0$: $\langle H_a v, v\rangle \geq 0$ para
todo $a$ y toda dirección $v$, o sea, todas las hessianas son
semidefinidas positivas.
\end{solution}

\begin{solution}{exo:b2:diffcalc:12}
\begin{enumerate}
  \item Por la desigualdad del valor medio
        (\cref{thm:b2:diffcalc:mvi}) aplicada a $g$: $\norm{g(x) -
        g(y)} \leq k\norm{x-y}$, luego
        \[
        \norm{f(x) - f(y)} \geq \norm{x - y} - \norm{g(x) -
        g(y)} \geq (1 - k)\norm{x - y} :
        \]
        $f$ es inyectiva y $f^{-1}$ (definida sobre la imagen) es
        $\frac{1}{1-k}$-lipschitziana.
  \item Fíjese $y$; $T(x) = y - g(x)$ cumple $\norm{T(x) - T(x')} =
        \norm{g(x') - g(x)} \leq k\norm{x - x'}$: es una contracción
        del espacio \omterm{def:b2:metric:complete}{completo} $\R^n$. Banach
        (\cref{thm:b2:metric:banach}) da un punto fijo $x^* = y -
        g(x^*)$, es decir, $f(x^*) = y$: $f$ es sobreyectiva.
  \item Así pues, $f$ es una biyección de $\R^n$ con inversa
        $\frac{1}{1-k}$-lipschitziana: un teorema de inversión
        global, donde la pequeñez de $\dd g$ en todas partes
        sustituye a la hipótesis de invertibilidad local del~\cref{thm:b2:diffcalc:inverse}.
\end{enumerate}
\end{solution}

\begin{solution}{pb:b2:diffcalc:1}
\textbf{1.} $\norm{ABx} \leq \vertiii A\norm{Bx} \leq \vertiii
A\vertiii B\norm x$: tómese el supremo sobre $\norm x = 1$. Los
productos de matrices, $\det$ y $\operatorname{tr}$ son funciones
polinómicas de las entradas, luego \omterm{def:b2:metric:continuity}{continuas} ($\mathcal M_n(\R) \simeq
\R^{n^2}$, con todas las \omterm{def:b2:nvs:equivalent}{normas equivalentes}:
\cref{thm:b2:nvs:finitedim}).

\textbf{2.} $\sum\vertiii{X^k} \leq \sum\vertiii X^k < \infty$: la
serie converge \omterm{def:b2:series:def}{absolutamente}, luego converge
(\cref{thm:b2:nvs:absoluteconvergence}). De $(I - X)\sum_{k\leq
N}X^k = I - X^{N+1} \to I$ se sigue que la suma es $(I - X)^{-1}$.
\omterm{def:b2:nvs:norm}{Norma}: $\leq \sum\vertiii X^k = \frac{1}{1 - \vertiii X}$; y
\[
(I - X)^{-1} - I - X = \sum_{k\geq2}X^k = X^2(I - X)^{-1},
\qquad
\vertiii{X^2(I-X)^{-1}} \leq
\frac{\vertiii X^2}{1 - \vertiii X} = O(\vertiii X^2).
\]

\textbf{3.} $A + H = A(I + A^{-1}H)$ con $\vertiii{A^{-1}H} \leq
\vertiii{A^{-1}}\vertiii H < 1$: invertible por la pregunta 2, de
modo que la bola \omterm{def:b2:metric:topology}{abierta} de radio $\vertiii{A^{-1}}^{-1}$ centrada en
$A$ está en $GL_n(\R)$. Densidad: $\det(A + \varepsilon I)$ es un
polinomio de grado $n$ en $\varepsilon$ con coeficiente principal
$1$: tiene un número finito de raíces, así que hay $\varepsilon_k \to
0$ con $A + \varepsilon_kI$ invertible, convergiendo a $A$.

\textbf{4.} Para $\vertiii H < \frac{1}{2\vertiii{A^{-1}}}$:
\[
(A + H)^{-1} - A^{-1}
= \bigl[(I + A^{-1}H)^{-1} - I\bigr]A^{-1},
\]
de \omterm{def:b2:nvs:norm}{norma} a lo sumo $\frac{\vertiii{A^{-1}H}}{1 -
\vertiii{A^{-1}H}}\,\vertiii{A^{-1}} \leq
2\vertiii{A^{-1}}^2\vertiii H \to 0$: $\Phi$ es \omterm{def:b2:metric:continuity}{continua} en todo $A
\in GL_n(\R)$.

\textbf{5.} $(A+H)^2 = A^2 + AH + HA + H^2$: la aplicación $H \mapsto
AH + HA$ es lineal y el error $H^2$ es $O(\vertiii H^2)$.
Desarrollando $(A + H)^k$ y ordenando por el número de factores $H$:
la parte lineal es $\sum_{i=0}^{k-1}A^iHA^{k-1-i}$, y los términos
con $\geq 2$ factores $H$ están acotados por $\binom k2$ productos de
\omterm{def:b2:nvs:norm}{norma} del orden de $\vertiii A^{k-2}\vertiii H^2$: $O(\vertiii
H^2)$. La suma no puede colapsarse a $kA^{k-1}H$ porque $H$ y $A$ no
tienen por qué conmutar: la suma es la derivada no conmutativa
correcta.

\textbf{6.} Para $H$ pequeño:
\[
(A+H)^{-1} = (I + A^{-1}H)^{-1}A^{-1}
= \bigl(I - A^{-1}H + O(\vertiii H^2)\bigr)A^{-1}
= A^{-1} - A^{-1}HA^{-1} + O(\vertiii H^2) :
\]
$\dd\Phi_A(H) = -A^{-1}HA^{-1}$, lineal en $H$. Para $n = 1$:
$\dd(1/a)(h) = -h/a^2$, la derivada de siempre.

\textbf{7.} Regla de la cadena a lo largo de la curva:
$\bigl(A(t)^{-1}\bigr)' = \dd\Phi_{A(t)}(A'(t)) =
-A(t)^{-1}A'(t)A(t)^{-1}$. En $A(t) = I + tB$, con $t = 0$: $(I +
tB)^{-1} = I - tB + O(t^2)$.

\textbf{8.} Por la pregunta 5 y la invariancia \omterm{def:b2:structures:generated}{cíclica} de la traza:
\[
\dd f_A(H) = \operatorname{tr}\Bigl(\sum_{i=0}^{k-1}
A^iHA^{k-1-i}\Bigr) = k\operatorname{tr}\bigl(A^{k-1}H\bigr) .
\]
Frente al producto de Frobenius, $\dd f_A(H) =
\operatorname{tr}\bigl((\nabla f)^{\mathsf T}H\bigr)$ exige $(\nabla
f)^{\mathsf T} = kA^{k-1}$: $\nabla f(A) = k\,(A^{k-1})^{\mathsf T}$.

\textbf{9.} $f(A + H) - f(A) = 2\operatorname{tr} (A^{\mathsf T}H) +
\operatorname{tr}(H^{\mathsf T}H)$: la \omterm{def:b2:diffcalc:differential}{diferencial} es $H \mapsto
2\operatorname{tr}(A^{\mathsf T}H) = 2\langle A, H\rangle$, luego
$\nabla f(A) = 2A$; y el término de segundo orden es exactamente
$\norm H_F^2$: la hessiana es el doble de la \omterm{def:b2:quadratic:def}{forma cuadrática}
identidad, definida positiva y constante, de modo que
$\norm\cdot_F^2$ es estrictamente convexa (aquí la fórmula de Taylor
es exacta).

\textbf{10.} Por multilinealidad en las columnas, $\det(I + H) =
\sum_{S\subseteq\{1,\dots,n\}}\det(M_S)$, donde $M_S$ tiene la
columna $h_j$ para $j \in S$ y $e_j$ en los demás casos. $S =
\varnothing$ da $1$; $S = \{j\}$ da el \omterm{def:b2:linalg:det}{determinante} de $I$ con la
columna $j$ sustituida por $h_j$, a saber, su entrada $j$-ésima
$h_{jj}$, que suman $\operatorname{tr}H$; cada término con $\abs S
\geq 2$ es un \omterm{def:b2:linalg:det}{determinante} con al menos dos columnas de tamaño
$O(\vertiii H)$, luego $O(\vertiii H^2)$ (las aplicaciones
multilineales sobre un espacio de dimensión finita son acotadas), y
hay un número finito de ellos. Por tanto, $\det(I + H) = 1 +
\operatorname{tr}H + O(\vertiii H^2)$: $\dd(\det)_I =
\operatorname{tr}$.

\textbf{11.} $\det(A + H) = \det A\,\det(I + A^{-1}H) = \det
A\,\bigl(1 + \operatorname{tr}(A^{-1}H) + O(\vertiii H^2)\bigr)$: la
\omterm{def:b2:diffcalc:differential}{diferencial} es $H \mapsto \det(A)\operatorname{tr}(A^{-1}H)$.

\textbf{12.} Desarrollo de Laplace por la fila $i$: $\det A = \sum_j
a_{ij}C_{ij}$, y los cofactores $C_{ij}$ no involucran la fila $i$:
$\frac{\partial\det}{\partial a_{ij}} = C_{ij}$. De ahí,
\[
\dd(\det)_A(H) = \sum_{i,j}C_{ij}h_{ij}
= \operatorname{tr}\bigl(\operatorname{com}(A)^{\mathsf T}
H\bigr)
= \operatorname{tr}\bigl(\operatorname{adj}(A)H\bigr),
\qquad
\nabla(\det)(A) = \operatorname{com}(A) .
\]
Para $A$ invertible, $\operatorname{adj}A = \det(A)A^{-1}$ recupera
la pregunta 11. A lo largo de una curva $C^1$, la regla de la cadena
se lee $(\det A(t))' =
\operatorname{tr}(\operatorname{adj}(A(t))\,A'(t))$: la fórmula de
Jacobi.

\textbf{13.} Con $A' = MA$ y $\operatorname{adj}(A)A = \det(A)I$:
\[
(\det A)' = \operatorname{tr}\bigl(\operatorname{adj}(A)MA
\bigr)
= \operatorname{tr}\bigl(A\operatorname{adj}(A)M\bigr)
= \det A\;\operatorname{tr}M
\]
(ciclicidad; y también $A\operatorname{adj}A = \det(A) I$). La EDO
lineal escalar $y' = \operatorname{tr}(M(t))\,y$ tiene la única
solución $y(t) = y(0)\exp\bigl(\int_0^t \operatorname{tr}M\bigr)$
(primer año): la fórmula de Liouville.

\textbf{14.} En $A \in SL_n(\R)$, tómese $H = \frac1nA$:
$\dd(\det)_A\bigl(\tfrac1nA\bigr) = \frac1n\det(A)
\operatorname{tr}(A^{-1}A) = \frac1n\cdot1\cdot n = 1 \neq 0$: la
\omterm{def:b2:diffcalc:differential}{diferencial} es una forma lineal no nula, luego sobreyectiva sobre
$\R$ en todo punto del conjunto de nivel: $SL_n(\R)$ es un conjunto
de nivel liso (de dimensión $n^2 - 1$).

\textbf{15.} $\sum_k\vertiii{A^k/k!} \leq \sum\vertiii A^k/k! =
\eu^{\vertiii A}$: convergencia absoluta (el argumento de
\omterm{def:b2:metric:complete}{completitud} de la pregunta 2), con \omterm{def:b2:funcseq:series}{convergencia normal} sobre toda
bola $\vertiii A \leq R$ (cota $R^k/k!$ independiente de $A$). Cada
suma parcial es \omterm{def:b2:metric:continuity}{continua} (polinómica); el límite uniforme sobre bolas
es \omterm{def:b2:metric:continuity}{continuo}: $A \mapsto \eu^A$ es \omterm{def:b2:metric:continuity}{continua}, con $\vertiii{\eu^A} \leq
\eu^{\vertiii A}$.

\textbf{16.} Ambas series convergen \omterm{def:b2:series:def}{absolutamente}, de modo que el
\omterm{thm:b2:series:fubini}{producto de Cauchy} es legítimo (\cref{thm:b2:series:fubini}):
\[
\eu^A\eu^B = \sum_{n\geq0}\frac{1}{n!}\sum_{k=0}^n\binom
nkA^kB^{n-k}
= \sum_{n\geq0}\frac{(A+B)^n}{n!} = \eu^{A+B},
\]
donde la identidad del binomio exige $AB = BA$. Con $B = -A$:
$\eu^A\eu^{-A} = \eu^0 = I$: toda $\eu^A \in GL_n(\R)$.

\textbf{17.} La serie $\sum t^kA^k/k!$ y su serie derivada $\sum
t^{k-1}A^k/(k-1)! = A\sum t^{k-1}A^{k-1}/(k-1)!$ convergen
\omterm{def:b2:funcseq:series}{normalmente} sobre todo segmento $\abs t \leq T$ (con cotas
$T^k\vertiii A^k/k!$): el teorema de derivación de series
(\cref{thm:b2:funcseq:seriestransfer}, aplicado entrada a entrada) da
$\frac{\dd}{\dd t}\eu^{tA} = A\eu^{tA}$; sacando factor $A$ por la
derecha se obtiene $\eu^{tA}A$. Iterando: es $C^\infty$.

\textbf{18.} $A(t) = \eu^{tA}$ cumple $A'(t) = A\,A(t)$: la fórmula
de Liouville (pregunta 13) con $M = A$ constante da $\det\eu^{tA} =
\eu^{t\operatorname{tr}A}$ (de valor $1$ en $t = 0$); en $t = 1$,
$\det\eu^A = \eu^{\operatorname{tr}A}$. Comprobaciones: $n = 1$ es la
propia exponencial; una $A$ nilpotente tiene $\operatorname{tr}A = 0$
y $\eu^A$ unipotente de \omterm{def:b2:linalg:det}{determinante} $1$; y $\operatorname{tr}A = 0$ da
$\det\eu^A = 1$: las matrices de traza nula van a parar a
$SL_n(\R)$.

\textbf{19.} $\eu^H - I - H = \sum_{k\geq2}H^k/k!$, de \omterm{def:b2:nvs:norm}{norma} $\leq
\vertiii H^2\eu^{\vertiii H} = O(\vertiii H^2)$: $\dd(\exp)_0 =
\mathrm{id}$, invertible. Además, $\exp$ es $C^1$: por la pregunta 5,
la \omterm{def:b2:diffcalc:differential}{diferencial} candidata $H \mapsto
\sum_k\frac1{k!}\sum_iA^iHA^{k-1-i}$ es una serie \omterm{def:b2:funcseq:series}{normalmente}
convergente de aplicaciones lineales que dependen continuamente de $A$
(con cotas $\vertiii A^{k-1}/(k-1)!$ sobre bolas), luego las
parciales existen y son \omterm{def:b2:metric:continuity}{continuas} (\cref{thm:b2:diffcalc:c1} y el
teorema de transferencia para series). El teorema de la función
inversa (\cref{thm:b2:diffcalc:inverse}) se aplica en $0$: $\exp$ es
un difeomorfismo $C^1$ de un entorno de $0$ sobre un entorno de $I$;
las matrices próximas a $I$ tienen logaritmo.

\textbf{20.} Para $S = PDP^{\mathsf T}$ \omterm{def:b2:quadratic:adjoint}{simétrica} (teorema espectral):
$\eu^S = P\eu^DP^{\mathsf T}$ es \omterm{def:b2:quadratic:adjoint}{simétrica} con \omterm{def:b2:reduction:eigen}{valores propios}
$\eu^{\lambda_i} > 0$: definida positiva. Sobreyectividad: una $Q =
P\operatorname{diag}(\mu_i)P^{\mathsf T}$ definida positiva ($\mu_i >
0$) es $\eu^S$ para $S = P\operatorname{diag}(\ln\mu_i)P^{\mathsf
T}$. Inyectividad: $\eu^S$ determina sus \omterm{def:b2:reduction:eigen}{subespacios propios}, que son
exactamente los de $S$ (sobre cada \omterm{def:b2:reduction:eigen}{subespacio propio} de $S$ para
$\lambda$, $\eu^S$ actúa como $\eu^\lambda$; y $\lambda$ distintos
dan $\eu^\lambda$ distintos), y tomar $\ln$ de los \omterm{def:b2:reduction:eigen}{valores propios}
recupera $S$. Así pues, $\exp$ es una biyección de las matrices
\omterm{def:b2:quadratic:adjoint}{simétricas} sobre las definidas positivas.

\textbf{21.} $F(A + H) = A^{\mathsf T}A + A^{\mathsf T}H + H^{\mathsf
T}A + H^{\mathsf T}H$: $\dd F_A(H) = A^{\mathsf T}H + H^{\mathsf T}A$
(con valores en $S_n$; error $O(\vertiii H^2)$). En $A \in O_n$ y
para $S \in S_n$, la elección $H = \frac12AS$ da
\[
A^{\mathsf T}\cdot\tfrac12AS + \tfrac12(AS)^{\mathsf T}A
= \tfrac12 S + \tfrac12 S^{\mathsf T} = S :
\]
sobreyectiva. $O_n = F^{-1}(I)$ es un conjunto de nivel liso de
dimensión $n^2 - \dim S_n = \frac{n(n-1)}2$.

\textbf{22.} Derivando $A(t)^{\mathsf T}A(t) = I$ en $t = 0$ (con
$A(0) = I$): $A'(0)^{\mathsf T} + A'(0) = 0$: antisimétrica.
Recíprocamente, para $K^{\mathsf T} = -K$: $(\eu^{tK})^{\mathsf
T}\eu^{tK} = \eu^{tK^{\mathsf T}} \eu^{tK} = \eu^{-tK}\eu^{tK} = I$
(traspóngase la serie término a término; los exponentes conmutan):
la curva permanece en $O_n$, con velocidad $K$ en $t = 0$. El espacio
tangente en $I$ es el de las matrices antisimétricas, de la dimensión
esperada $\frac{n(n-1)}2$.

\textbf{23.} $\operatorname{tr}K = 0$ para $K$ antisimétrica, luego
$\det\eu^K = \eu^0 = 1$ (pregunta 18): la exponencial aterriza en
$SO_n$. Para $n = 2$: $J^2 = -I$, luego $J^{2m} = (-1)^mI$ y
$J^{2m+1} = (-1)^mJ$, y
\[
\eu^{\theta J}
= \Bigl(\sum_m\frac{(-1)^m\theta^{2m}}{(2m)!}\Bigr)I
+ \Bigl(\sum_m\frac{(-1)^m\theta^{2m+1}}{(2m+1)!}\Bigr)J
= \cos\theta\,I + \sin\theta\,J ,
\]
la rotación de ángulo $\theta$: las series del seno y del coseno
viven dentro de la \omterm{ex:b2:nvs:matrixexp}{exponencial de matrices}.

\textbf{24.} Para $A, B$ invertibles: $\operatorname{adj}(AB) =
\det(AB)(AB)^{-1} = \det(B)\det(A)B^{-1}A^{-1} =
\operatorname{adj}(B)\operatorname{adj}(A)$. Ambos miembros de la
identidad son aplicaciones polinómicas (luego \omterm{def:b2:metric:continuity}{continuas}) de las
entradas de $(A, B)$; coinciden sobre el subconjunto denso
$GL_n\times GL_n$ de $\mathcal M_n\times\mathcal M_n$ (pregunta 3:
aproxímese cada factor), luego coinciden en todas partes.

\textbf{25.} (i) La parte I funcionó con la \omterm{def:b2:metric:complete}{completitud} de los espacios
normados de dimensión finita (las series \omterm{def:b2:series:def}{absolutamente} convergentes
convergen); la parte IV, con lo mismo más el teorema espectral para
la pregunta 20; y el logaritmo local de la parte V, con el teorema de
la función inversa. (ii) El capítulo de ecuaciones diferenciales se
apoya en la fórmula de Liouville (pregunta 13) para el wronskiano de
los sistemas lineales, y en $\frac{\dd}{\dd t}\eu^{tA} = A\eu^{tA}$
(pregunta 17), que es el enunciado de que $\eu^{tA}$ resuelve $X' =
AX$. (iii) $\dd(\det)_I = \operatorname{tr}$ dice que la traza es la
tasa infinitesimal de cambio de volumen, y $\det\eu^A =
\eu^{\operatorname{tr}A}$ integra globalmente ese enunciado. (iv) El
volumen del tercer año pone nombre a las estructuras: $O_n$ y
$SL_n(\R)$ son grupos de Lie, sus espacios tangentes en $I$ (las
matrices antisimétricas y las de traza nula) son álgebras de Lie, y
$\exp$ es el puente entre unos y otras.
\end{solution}
